\documentclass[10pt,a4paper]{amsart}
\usepackage[margin=19mm]{geometry}
\usepackage{amsmath,amssymb,mathtools}
\usepackage[scr=rsfs,cal=euler]{mathalfa}
\usepackage{xfrac}
\usepackage[unicode,hidelinks,bookmarks=false]{hyperref}
\newcommand{\ie}{i.e.\ }
\newcommand{\cf}{cf.\ }
\newcommand{\resp}{respectively\ }
\newcommand{\Subsection}{\subsection}
\providecommand{\nobreakdash}{\nobreak\hbox{-}}
\setlength{\emergencystretch}{2em}
\multlinegap0pt
\usepackage{booktabs}
\usepackage{natbib}
\usepackage{xcolor}
\usepackage{tikz}
\usepackage{mathtools}
\usepackage{siunitx}
\newtheorem{corollary}{Corollary}
\title{Metric geometry for ranking-based voting: Tools for learning electoral structure}
\author{Moon Duchin, Kristopher Tapp}
\date{Source: arXiv:2602.10293v1 (CC BY 4.0)}
\begin{document}
\maketitle

\noindent\emph{Source and licence.} Metric geometry for ranking-based voting: Tools for learning electoral structure. Version arXiv:2602.10293v1 (CC BY 4.0), distributed by arXiv under the Creative Commons Attribution 4.0 International licence, http://creativecommons.org/licenses/by/4.0/, as recorded in the arXiv licence metadata for this exact version and shown on the abstract page https://arxiv.org/abs/2602.10293v1. Full licence notice: \texttt{licenses/arxiv-2602.10293-CC-BY-4.0.txt}.\par\smallskip

\noindent\textit{Editorial scope.} Counted unit: the article's corollary P:my\_diaconis (source0.tex lines 444--447). The article's complete original proof of it is delivered with it (source0.tex lines 448--455, one \texttt{proof} environment of 8 lines). No mathematics is written, repaired or re-derived: the statement and the proof are the article's own lines, in the article's own notation, and no retained line is substituted or re-flowed.  Retained uncounted as the source-local context the proof needs: lines 427--429.  Nothing was omitted from any retained span, so the package carries no omission record.  No retained line contains a citation, so the wrapper carries no bibliography and no bibliography entry.  No retained line contains a figure environment, an image-inclusion command or drawing code, so no image is delivered and the package declares no asset dependency.  Editorial formatting only, in addition: an AMS \texttt{amsart} preamble that restates the article's own declarations and macros used by the retained lines, namely the environments \texttt{corollary} on separate counters, together with the standard packages those lines need.  The retained environments print in this document as Corollary 1 in order of appearance, whereas the article prints the same items with the numbering of its own full document.  The complete native source of the article as distributed in the arXiv e-print for this exact version is bundled unchanged as \texttt{sources/194383/source0.tex}.
\begin{corollary}\label{C:hihi}
$d_B\leq d_H < 2d_B$.
\end{corollary}

\begin{corollary}\label{P:my_diaconis} For any pair $\sigma,\tau\in\Omega_n$ of distinct ballots,
$\displaystyle\frac{d_H(\sigma,\tau)}{d_B(\sigma,\tau)}\in[1,2)$,
and these bounds are sharp.
\end{corollary}

\begin{proof}
The bound is just a rephrasing of Corollary~\ref{C:hihi}, so it remains to establish optimality.
If $\sigma\in\Omega_m$ is a complete ballot, let 
$\hat\sigma$ be its reversal, with candidates listed in opposite order.  We have $d_H(\sigma,\hat\sigma)=\binom{m}{2}$ because every pair will have a strong disagreement.  
On the other hand, $d_B(\sigma,\hat\sigma)= (m-1)(m+1)$, so comparing these gives a sequence with 
$\displaystyle\frac{d_H(\sigma,\hat\sigma)}{d_B(\sigma,\hat\sigma)}
\to 2$.  For the other bound, just note that any pair of ballots connected by a single neighbor swaps realizes $d_H=d_B$.
\end{proof}

\end{document}
